# Profil E1 — Prise de notes brute (Manager, 10 ans)

**Code :** E1  
**Fonction :** Manager  
**Ancienneté :** 10 ans  
**Contexte :** Banque, assurances, institutions européennes (EIF). Chef de projet, livrables techniques.

---

## T0 — Cadrage

**Q : Parcours, comment arrivé ici ?**
- 10 ans. Commencé analyste.
- Banque, assurances. Instit européennes aussi, EIF. Livrables techniques.
- Encadre 1-2 juniors maintenant. Avant 3-4.
- Beaucoup pilotage. Chef de projet. PPT, suivis, revues.
- À mon niveau, passe + temps à relire qu'à écrire.

---

## T1 — Usage réel

**Q : Dernière fois utilisé outil génératif ?**
- 2-3 semaines. Curiosité. Junior parlait d'un prompt pour notes réunion.
- Testé sur CR que j'avais fait. Résultat correct, mais pas + que ce que j'avais.
- À mon niveau, sais ce que veux. Pas besoin d'expliquer 2h à machine. Plus rapide seul.

**Q : Sur quelles tâches au quotidien ?**
- Pas vraiment quotidien. De temps en temps Copilot. Sujets moins maîtrisés. Rare.
- Préfère méthodes : templates, bases données anciens livrables, connaissance clients.
- IA apporte pas grand-chose.
- Juniors l'utilisent tout le temps. Eux ont besoin de béquille.

**Q : Comment appris ?**
- Session sensibilisation y a qqs mois. Pas creusé. Pas de bloc-notes prompts, pas veille.
- Sais à peu près ce que ça fait. Jamais dit "faut maîtriser".
- À mon niveau, résultat final compte. Pas processus.

**Q : Temps gagné mis sur quoi ?**
- Pas vraiment, car utilise pas.
- Mais gain de temps à mettre sur relecture/vérification. Ce qui sort IA demande contrôle.
- Au final, est-ce qu'on gagne vraiment du temps ? Pas sûr.
- Préfère faire une fois bien que corriger approximatif.

**Q : Situation où plus d'effort que de le faire soi-même ?**
- Il y a qqs mois. Livrable client important. Junior avait utilisé IA sur tout.
- Pas plongé dans projet avant. Confiance relative. Deadline.
- Reçois doc. Incohérences partout. Chiffres faux. Réfs à textes inexistants. Formulations justes mais sens trahi.
- Refait tout en catastrophe, avec junior, une nuit.
- Passé + temps corriger IA que produire nous-mêmes.
- Depuis, très prudent. Demande systématiquement : "qu'est-ce qui est de toi, qu'est-ce qui est de l'IA ?"

**Q : Vous en parlez entre vous ?**
- Entre managers, oui. On partage. Consensus : IA là, faut composer, œil critique.
- Tous eu histoires livrables problématiques. Certains refaits entièrement.
- À mon niveau, changement métier. Vérifie plus seulement fond/forme. Vérifie aussi ce que IA a apporté. Compétence nouvelle à développer.

---

## T2 — Apprentissage

**Q : Comment on apprend le métier aujourd'hui ?**
- Quand j'ai commencé : faisait, se trompait, se faisait reprendre.
- Aujourd'hui juniors arrivent avec livrables très propres très vite. Mais profondeur ?
- IA donne première version, retravaillent, mais pas chemin réflexion initiale.
- Jeunes très bons forme, mais mal justifier fond. Apprentissage différent. Pas sûr meilleur.

**Q : Par quelles tâches on apprend vraiment ?**
- Choix difficiles. Décider reco, orientation stratégique, pertinence client. IA peut pas.
- Répétitives, rédaction, mise en forme : IA peut aider. Peut-être bonne chose.
- Permet juniors passer + temps sur valeur ajoutée.
- Mais faut qu'ils aient fait tâches ingrates au moins 1-2x à la main. Pour comprendre. Sans, sauront pas juger ce que IA sort.

**Q : Comment vérifiez ce que produit junior qui utilise IA ?**
- Plus seulement contenu final. Processus.
- Demande : "comment arrivé ? Qu'est-ce fait toi ? Qu'est-ce IA ? Vérifié sources ?"
- Veux savoir ce qui vaut quelque chose ou pas, d'après lui.
- IA peut être brillant ou catastrophique. Junior doit savoir faire diff.
- Devenu partie importante rôle manager : former juniors à vérification critique.

**Q : Relecture senior, qu'est-ce se transmet ?**
- Jugement. Pourquoi telle formulation meilleure ? Pourquoi reco pertinente ? Pourquoi chiffre suspect ?
- Relecture = transmission compétences, pas correction technique.
- Avec IA, transmission + importante. Juniors peuvent illusion outil fait travail.
- Mon rôle = montrer que non, jugement reste leur responsabilité.

---

## T3 — Client et valeur

**Q : Confiance client repose sur quoi ?**
- Solidité. Banque/assurance, pas approximatif.
- Client doit sentir vérifié, solide, pris responsabilités.
- Si IA a fait erreurs et pas vues, confiance rompue. Se reconstruit pas facilement.

**Q : Sujet IA déjà abordé avec client ?**
- Oui. Clients demandent si on utilise IA.
- Répond : technologies aide décision, mais analyse finale/recommandations entièrement équipes.
- C'est vrai. Mais jamais "on génère slides avec Copilot". Mal perçu. Client se demanderait ce qu'il paie.

**Q : Diriez-vous au client qu'une partie est IA ?**
- Non. Contre-productif. Client paie résultat, pas processus.
- Si résultat bon, peu importe comment.
- Si mauvais, pas cacher derrière IA. Responsabilité sur nous.

**Q : Livrable moins crédible ?**
- Erreurs factuelles. Junior IA sans vérifier, coquilles grossières.
- Client pas dit direct, mais méfiance après. Mission suivante a demandé autre consultant.
- Depuis, intransigeant relecture. À mon niveau, peux pas laisser passer erreur.

**Q : Dans 10 ans, client paiera pour quoi ?**
- Jugement et responsabilité. IA fera +, mais pas prendre risque à place consultant.
- Banque/assurance, prise risque centrale. Client paie quelqu'un qui dit "je recommande, j'assume".
- IA peut pas. Valeur cabinet = jugement, pas production.

---

## T4 — Gouvernance

**Q : Règles sur usage ?**
- Politique. Copilot officiel. Règles confidentialité.
- À mon niveau, règles un peu légères. Flou outils non-officiels. Pas processus vérification spécifique.
- Repose responsabilité individuelle. Bien, mais pas suffisant vu risques.

**Q : Avant livrable client ?**
- Relecture, moi ou manager. Jusqu'à récent, pas vérifiait spécifiquement IA.
- Depuis incident, vérifie. Demande sources, justification choix.
- Réflexe sécurité. Managers doivent développer compétence vérification spécifique. IA change nature erreurs.

**Q : Déjà rencontré contenu faux IA ?**
- Oui, livrable entier. Chiffres faux, références inventées, structure pas tenait.
- Heureusement refait avant client. Nuit à reprendre. Livré qqs heures retard.
- Junior leçon, moi aussi. Changé supervision. Pas confiance aveugle. Tout contrôler.

**Q : Règle idéale ?**
- "IA outil assistance, pas production. Chaque livrable justifiable sans recours outil."
- Traçabilité : savoir ce qui a été produit par IA et pas. Pas contrôler, pouvoir vérifier si problème.
- À mon niveau, minimum éviter catastrophes.

---

## T5 — Clôture

**Q : À la place de la direction, première décision ?**
- Système validation renforcé pour livrables utilisant IA. Pas interdiction, marquage.
- Mention "préparé avec assistance IA" déclenche relecture renforcée.
- Former managers à relecture spécifique. Erreurs IA différentes humaines. + subtiles, + difficiles détecter.

**Q : Aspect important non abordé ?**
- Pression commerciale. Clients demandent + IA. Cabinet doit s'adapter. Managers doivent montrer exemple.
- Contradictoire : on dit maîtriser IA, mais pas formation, pas temps. À mon niveau, garder compétences techniques et apprendre nouveaux outils. Pression supplémentaire.

**Q : Point de vue utile à recueillir ?**
- Directeur ou associé. Voir comment ils voient évolution métier, intégration IA stratégie cabinet.
- Moi vois risques livrables. Eux risques marque/réputation. Pression clients qui demandent solutions IA. Regard très différent.